\section{Society5.0：技術融合と人間中心の目的}
\label{theory:society:five}
狩猟社会、農耕社会、工業社会、情報社会に続く第五の社会という整理を読む。サイバー空間と現実空間の融合によって、経済的発展と社会的課題の解決の両立を目指す。中心に置くのは人間であり、技術の導入量だけを目的にしない\parencite{eth54_society2021}。

2016年の提案と2021年の計画を分ける。2021年には持続可能性、強靱性、安全・安心、多様な幸福への説明が具体化される。人文・社会科学も含む知識をつなぎ、倫理・法・社会への影響を検討する。過去問の穴埋めでは、社会段階、社会的課題、人間中心、持続可能性を別の空欄として確かめる。後年の政策用語を過去の指定語へ無条件に置き換えない。
